% Pursuing Stacks (cotes 134-2 to 134-8): archive cote 134-4, site page 92 (PDF page 93), typescript page 179.
% typescript. Draft compared: ps3.tex:4736-4787.
% AI-transcribed by Claude Opus 5.5 (claude-opus-5-5, adaptive thinking, effort high), Message Batches API, batch ps2 (msgbatch_01NVwr1Yuo8qJf5rykeGBXPK), 2026-11-12; image review recorded in records/CHECK_DECISIONS_134_3_4.tsv.
\scanpage{134-4/92}
\mspage{179}
objects of $A^{\wedge}$ contractible for the homotopy structure $h_{\underline{W}_A}$ associated to
$\underline{W}_A$ - or equivalently, $i^{*}$ is a morphism of homotopy structures(cf. \inserted{defi-}%
\inserted{nition on} page 134).

(vi) T\over[typed]{h}{H}e functor $i^{*}$ transforms the projection $\hand{\Delta}_1 \hand{\to} e$ into a\xout{n}
\xout{element of $\underline{W}$} map in $\underline{W}$ , or equivalently (as $A$ is $\underline{W}$-aspheric) $i^{*}(\hand{\Delta}_1)$
is $\underline{W}$-aspheric.

We have the trivial implications
\[ (i)\hand{\Rightarrow}(ii)\hand{\Rightarrow}(iii)\hand{\Rightarrow}(iv)\hand{\Rightarrow}
\underset{\text{\deleted{\unsure{(vi)}}}}{\overset{\substack{\text{\deleted{\xout{\unsure{(v)}}}}\\\text{\deleted{\hand{\unsure{(vi)}}}}}}{(vi)}}
\hand{\Longleftarrow}\ \underset{\text{\deleted{\illegible{1}}}}{\text{\inserted[typed]{(v)}}}\quad , \]
where the implication (v)$\hand{\Rightarrow}$(vi) is in fact an equivalence, due to the
fact that the contractibility structure on (Cat) is defined in terms of
$\hand{\Delta}_1$ as a generating \xout{homotopy} contractor, and that the assumption $i^{*}(\hand{\Delta}_1)$
$\underline{W}$-aspheric implies already that it is $\underline{W}_A$ -aspheric \inserted[typed]{over $e_{A^{\wedge}}$} (because $A$ is totally
$\underline{W}_A$ -aspheric), and hence contractible as it is a contractor and $h_{\underline{W}_A}$ is
defined in terms of \inserted[typed]{``weak homotopy} \deleted{\xout{\illegible{1}}} intervals'' which are $\underline{W}_A$-aspheric over $e_{A^{\wedge}}$. Thus
the only delicate implication is (vi)$\hand{\Rightarrow}$(i), which however follows from
\hand{(}theorem 1\inserted{, (ii) $\Rightarrow$ (i))}.

We got the longed-for ``key result'' \deleted{after all} \inserted{in the end} !

\margin{92}
